\documentclass[12pt, a4paper]{article}

% --- UNIVERSAL PREAMBLE BLOCK ---
\usepackage[a4paper, top=2.5cm, bottom=2.5cm, left=2.5cm, right=2.5cm]{geometry}
\usepackage{fontspec}

% تنظیمات زبان و جهت‌چینش برای موتور LuaLaTeX
% استفاده از babel به جای xepersian برای سازگاری با محیط
\usepackage[bidi=basic, layout=counter, provide=*]{babel}

\babelprovide[import, main, mapdigits]{persian}
\babelprovide[import]{english}

% تنظیم فونت‌ها (استفاده از فونت‌های استاندارد Noto موجود در سیستم)
\babelfont{rm}{Vazirmatn}
\babelfont[persian]{rm}{Vazirmatn}

% اصلاح لیست‌ها برای زبان فارسی
\usepackage{enumitem}
\setlist[itemize]{label=-}

% بسته‌های مورد نیاز برای جدول و ریاضی
\usepackage{amsmath}
\usepackage{booktabs}
\usepackage{array}
\usepackage{float}

\title{\vspace{-2cm} \textbf{پیش‌گزارش آزمایش هشتم: طراحی تقسیم‌کننده دودویی}}
\author{نام دانشجو: محمدرضا حسن‌زاده  \\ شماره دانشجویی: 40226046}
\date{تاریخ تحویل: \today}

\begin{document}

\maketitle

\section*{۱. هدف آزمایش}
هدف از این آزمایش، طراحی و پیاده‌سازی سخت‌افزاری یک مدار تقسیم‌کننده دودویی (Binary Divider) با استفاده از زبان توصیف سخت‌افزار (VHDL) است. در این مدار، یک عدد ۸ بیتی (به عنوان مقسوم) بر یک عدد ۴ بیتی (به عنوان مقسوم‌علیه) تقسیم می‌شود و حاصل‌ضرب به صورت خارج‌قسمت و باقی‌مانده در خروجی نمایش داده می‌شود.

\section*{۲. شرح الگوریتم (Restoring Division)}
برای پیاده‌سازی این مدار از روش «تقسیم با بازگرداندن باقی‌مانده» استفاده می‌شود. روند کلی این الگوریتم مشابه تقسیم دستی اعداد است. 
در این روش از سه ثبات اصلی استفاده می‌شود:
\begin{itemize}
    \item \textbf{ثبات $A (Accumulator)$:} برای نگهداری باقی‌مانده (مقدار اولیه صفر).
    \item \textbf{ثبات $Q (Quotient)$:} ابتدا حاوی مقسوم است و در انتها خارج‌قسمت را نگه می‌دارد.
    \item \textbf{ثبات $M (Divisor)$:} حاوی مقسوم‌علیه.
\end{itemize}

\textbf{مراحل انجام کار (تکرار به تعداد بیت‌ها):}
\begin{enumerate}
    \item شیفت منطقی ترکیب ثبات‌های $A$ و $Q$ به سمت چپ (Shift Left AQ).
    \item تفریق مقسوم‌علیه از ثبات $A$ ($A \leftarrow A - M$).
    \item بررسی بیت علامت $A$:
    \begin{itemize}
        \item اگر $A$ منفی شد: عمل بازگرداندن (Restore) انجام می‌شود ($A \leftarrow A + M$) و بیت کم‌ارزش $Q$ صفر می‌شود.
        \item اگر $A$ مثبت ماند: نیازی به بازگرداندن نیست و بیت کم‌ارزش $Q$ یک می‌شود.
    \end{itemize}
\end{enumerate}

\section*{۳. مثال عددی}
برای بررسی صحت عملکرد الگوریتم، تقسیم عدد ۵ ($0101_2$) بر ۳ ($0011_2$) را در نظر می‌گیریم. \\
\textbf{فرض‌ها:} $N=4$ (تعداد مراحل)، $Q=0101$، $M=0011$ و $A=0000$.

\begin{table}[H]
\centering
\caption{جدول ردیابی مراحل تقسیم ۵ بر ۳}
\vspace{0.3cm}
\begin{tabular}{|c|l|c|c|l|}
\hline
\textbf{مرحله} & \textbf{عملیات} & \textbf{ثبات A} & \textbf{ثبات Q} & \textbf{توضیحات} \\
\hline
شروع & مقداردهی اولیه & $0000$ & $0101$ & وضعیت اولیه \\
\hline
 & Shift Left AQ & $0000$ & $1010$ & شیفت به چپ \\
1 & Subtract $(A-M)$ & $1101$ & $1010$ & $0 - 3 = -3$ (منفی) \\
 & Check \& Restore & $\mathbf{0000}$ & $1010$ & بازگردانی A، بیت $Q_0=0$ \\
\hline
 & Shift Left AQ & $0001$ & $0100$ & شیفت به چپ \\
2 & Subtract $(A-M)$ & $1110$ & $0100$ & $1 - 3 = -2$ (منفی) \\
 & Check \& Restore & $\mathbf{0001}$ & $0100$ & بازگردانی A، بیت $Q_0=0$ \\
\hline
 & Shift Left AQ & $0010$ & $1000$ & شیفت به چپ \\
3 & Subtract $(A-M)$ & $1111$ & $1000$ & $2 - 3 = -1$ (منفی) \\
 & Check \& Restore & $\mathbf{0010}$ & $1000$ & بازگردانی A، بیت $Q_0=0$ \\
\hline
 & Shift Left AQ & $0101$ & $0000$ & شیفت به چپ \\
4 & Subtract $(A-M)$ & $0010$ & $0000$ & $5 - 3 = +2$ (مثبت) \\
 & Set Q Bit & $\mathbf{0010}$ & $\mathbf{0001}$ & عدم نیاز به بازگردانی، $Q_0=1$ \\
\hline
\end{tabular}
\end{table}

\textbf{نتیجه نهایی:}
پس از ۴ مرحله، مقدار ثبات $Q$ برابر با $0001$ (خارج‌قسمت = ۱) و مقدار ثبات $A$ برابر با $0010$ (باقی‌مانده = ۲) می‌باشد.

\end{document}